\documentclass[../preamble]{subfiles}
\begin{document}

\section{Monoidal Categories}
\label{sec:monoidal-category}
\concept{monoidal-category}{Monoidal Category}
\depends{natural-isomorphism}
\depends{bifunctor}
\implements{haskell}{Cat.Monoidal}
\implements{agda}{Cat.Monoidal}
\implements{lisp}{monoidal-category}

A \newterm{monoidal category} equips a category with a notion of ``multiplication''
of objects and morphisms,
coherently associative and unital up to natural isomorphism
(\nlabref{monoidal+category};
Mac Lane, \emph{Categories for the Working Mathematician}, Chapter VII).

\begin{definition}[Monoidal category]
  \label{def:monoidal-category}
  A \newterm{monoidal category}
  $\newmath{(\Category{C}, \otimes, I, \alpha, \lambda, \rho)}$
  consists of the following data:
  \begin{enumerate}
    \item A category $\Category{C}$.

    \item A \newterm{tensor product} bifunctor
      (Definition~\ref{def:bifunctor})
      \begin{equation}
        \label{eq:tensor-product}
        \newmath{\otimes \colon \Category{C} \times \Category{C} \lto \Category{C}}.
      \end{equation}
      We write $\newmath{A \otimes B}$ for $\otimes(A, B)$
      on objects,
      and $\newmath{f \otimes g}$ for $\operatorname{bimap}(f, g)$
      on morphisms.

    \item A \newterm{unit object}
      $\newmath{I \in \Ob(\Category{C})}$.

    \item An \newterm{associator},
      a natural isomorphism
      (Definition~\ref{def:nat-iso-componentwise})
      with components
      \begin{equation}
        \label{eq:associator}
        \newmath{\alpha_{A,B,C}
          \colon (A \otimes B) \otimes C
          \xrightarrow{\;\sim\;} A \otimes (B \otimes C)}
      \end{equation}
      natural in $A$, $B$, and $C$.

    \item A \newterm{left unitor},
      a natural isomorphism with components
      \begin{equation}
        \label{eq:left-unitor}
        \newmath{\lambda_A \colon I \otimes A \xrightarrow{\;\sim\;} A}
      \end{equation}
      natural in $A$.

    \item A \newterm{right unitor},
      a natural isomorphism with components
      \begin{equation}
        \label{eq:right-unitor}
        \newmath{\rho_A \colon A \otimes I \xrightarrow{\;\sim\;} A}
      \end{equation}
      natural in $A$.
  \end{enumerate}

  These data are subject to the pentagon axiom
  (Definition~\ref{def:pentagon})
  and the triangle axiom
  (Definition~\ref{def:triangle}).
\end{definition}

\subsection{Coherence Axioms}

\begin{definition}[Pentagon axiom]
  \label{def:pentagon}
  \axiom{monoidal-category}{pentagon}
  The \newterm{pentagon axiom} requires that
  for all objects $A, B, C, D \in \Ob(\Category{C})$,
  the following diagram commutes:
  \begin{equation}
    \label{eq:pentagon}
    \begin{tikzcd}[column sep=small]
      & ((A \otimes B) \otimes C) \otimes D
        \arrow[dl, "\alpha_{A,B,C} \otimes \id_D"']
        \arrow[dr, "\alpha_{A \otimes B, C, D}"]
      & \\
      (A \otimes (B \otimes C)) \otimes D
        \arrow[d, "\alpha_{A, B \otimes C, D}"']
      & &
      (A \otimes B) \otimes (C \otimes D)
        \arrow[d, "\alpha_{A, B, C \otimes D}"]
      \\
      A \otimes ((B \otimes C) \otimes D)
        \arrow[rr, "\id_A \otimes \alpha_{B,C,D}"']
      & &
      A \otimes (B \otimes (C \otimes D))
    \end{tikzcd}
  \end{equation}
  That is,
  the two composite paths from
  $((A \otimes B) \otimes C) \otimes D$
  to
  $A \otimes (B \otimes (C \otimes D))$
  are equal.
\end{definition}

\begin{definition}[Triangle axiom]
  \label{def:triangle}
  \axiom{monoidal-category}{triangle}
  The \newterm{triangle axiom} requires that
  for all objects $A, B \in \Ob(\Category{C})$,
  the following diagram commutes:
  \begin{equation}
    \label{eq:triangle}
    \begin{tikzcd}
      (A \otimes I) \otimes B
        \arrow[rr, "\alpha_{A,I,B}"]
        \arrow[dr, "\rho_A \otimes \id_B"']
      & &
      A \otimes (I \otimes B)
        \arrow[dl, "\id_A \otimes \lambda_B"]
      \\
      & A \otimes B &
    \end{tikzcd}
  \end{equation}
  That is,
  $(\id_A \otimes \lambda_B) \circ \alpha_{A,I,B}
    = \rho_A \otimes \id_B$.
\end{definition}

\subsection{Remarks}

\begin{remark}[Coherence]
  \label{rem:coherence}
  Mac Lane's \newterm{coherence theorem} for monoidal categories states that
  every diagram whose edges are built from
  $\alpha$, $\lambda$, $\rho$, their inverses,
  identities, and tensor products of these
  commutes.
  In particular,
  the pentagon and triangle axioms suffice to ensure coherence
  of all re-bracketing and unit-insertion paths.
  See Mac Lane, \emph{CWM} VII.2.
\end{remark}

\begin{remark}[Strictness]
  \label{rem:strictness}
  This definition is \emph{not} that of a strict monoidal category.
  The associator and unitors are natural isomorphisms,
  not identities.
  A \newterm{strict monoidal category} is one where
  $\alpha$, $\lambda$, and $\rho$ are all identity natural transformations,
  so that $(A \otimes B) \otimes C = A \otimes (B \otimes C)$
  and $I \otimes A = A = A \otimes I$ on the nose.
  By Mac Lane's coherence theorem,
  every monoidal category is monoidally equivalent to a strict one.
\end{remark}

\begin{remark}[Naturality spelled out]
  \label{rem:naturality}
  The naturality of $\alpha$ means that
  for all morphisms
  $f \colon A \lto A'$,
  $g \colon B \lto B'$,
  $h \colon C \lto C'$,
  the following square commutes:
  \begin{equation}
    \label{eq:associator-naturality}
    \begin{tikzcd}
      (A \otimes B) \otimes C
        \arrow[r, "\alpha_{A,B,C}"]
        \arrow[d, "(f \otimes g) \otimes h"']
      & A \otimes (B \otimes C)
        \arrow[d, "f \otimes (g \otimes h)"]
      \\
      (A' \otimes B') \otimes C'
        \arrow[r, "\alpha_{A',B',C'}"']
      & A' \otimes (B' \otimes C')
    \end{tikzcd}
  \end{equation}
  Analogous naturality squares hold for $\lambda$ and $\rho$.
\end{remark}

\end{document}
